\documentclass{article}
\usepackage[utf8]{inputenc}
\usepackage[ngerman]{babel}

\title{Anforderungen Master-Projektseminar: Echtzeitsysteme (SS26)}
\author{Danielou Mounsande}
\date{April 2026}

\begin{document}

\maketitle

\section{Generelle Vorgaben}
\begin{itemize}
    \item \textbf{Programmiersprache:} Ausschließlich C.
    \item \textbf{IDE:} STM32 Cube MX und STM32 Cube IDE.
    \item \textbf{Git:} Nur der \texttt{main} Branch wird bewertet. Saubere Dokumentation der Eigenleistung in Commit-Messages und Code (Namen an Funktionen/Dateien).
    \item \textbf{Instrumentierung:} Segger SystemView (via RTT).
    \item \textbf{Verifikation:} TeSSLa Spezifikationssprache.
\end{itemize}

\section{RTOS Spezifikationen}
\begin{itemize}
    \item \textbf{Tasks:} Fest zur Laufzeit, feste Priorität.
    \item \textbf{Scheduling:} Präemptives prioritätsbasiertes Round-Robin.
    \item \textbf{ISRs:} Rückkehr zum ursprünglichen Code (kein automatischer Kernel-Wechsel).
    \item \textbf{Hardware-Konfiguration:} FPU Deaktivieren, Soft-FP aktivieren. Nur MSP (Main Stack Pointer) nutzen.
    \item \textbf{Datenstrukturen:} Nutzung von \texttt{tcb.h} (vorgegeben).
\end{itemize}

\section{Kernel-Objekte}
\begin{itemize}
    \item \textbf{Semaphoren/Mutexes:} Acquire \& Release (Blocking \& Non-Blocking).
    \item \textbf{Delay:} Blocking (While-Loop) und Non-Blocking (Blocked-State).
    \item \textbf{Message Queues:} Send, Read, IsEmpty, IsFull (Blocking \& Non-Blocking).
\end{itemize}

\section{Referenzen}
\begin{itemize}
    \item Segger SystemView: \texttt{https://github.com/SEGGERMicro/SystemView}
    \item STM32L4 HAL: \texttt{UM1884}
    \item Cortex-M4 Programming Manual: \texttt{PM0214}
    \item Qing Li: \textit{Real-Time Concepts for Embedded Systems}
\end{itemize}

\end{document}
